\begin{papercard}{chanenson2025}
\cf{Stages} \stagetag{L} (prompted against fine-tuned LLMs for span annotation of policy sentences), \stagetag{C} (contextual-integrity parameters as a coding scheme), \stagetag{H} (information flows as the unit of harm reasoning). PoPETs 2025 paper by the group behind the CI crowdsourcing work of Shvartzshnaider et al.; added from the 23 October 2026 gap search because it is the clearest test of prompting on a theory-laden span task.

\cf{Task} Contextual integrity describes a flow by five parameters \pc{chanenson2025}{3}{1) the sender of the information, 2) the recipient of the information, 3) the subject of the information, 4) the information content or attribute, and 5) the transmission principle}; the GKC extension splits the transmission principle into aims, conditions, modalities and consequences. Annotation means marking these spans in a policy sentence, e.g.\ ``we'' as recipient and ``contact information'' as attribute \pc{chanenson2025}{3}{would label “we” as a recipient, “contact information” as an attribute, “you” as the sender}. Sentences were the input unit \pc{chanenson2025}{5}{We used sentences as the atomic unit for model input because sentence divisions are natural delimiters}.

\cf{Design} Ground truth: 21,588 expert GKC-CI annotations from 16 policies; 50 fine-tuned models from five families, open-source and proprietary \pc{chanenson2025}{1}{We fine-tune 50 open-source and proprietary models on 21,588 ground truth GKC-CI annotations from 16 privacy policies}. Baselines: an RNN and prompted GPT-4, GPT-4 Turbo and GPT-3.5 Turbo with 1, 3 or 5 random examples. Scoring is strict: a correct annotation is an exact string match with the expert span \pc{chanenson2025}{2}{We consider an accurate annotation to be an exact string match of the human annotator’s text}.

\cf{Results} Prompting failed: the best prompted model stayed below 20\,\% test accuracy \pc{chanenson2025}{7}{The prompted non-fine-tuned LLMs also performed poorly, with the best such model having a test set accuracy of <20\%}, and the authors conclude that off-the-shelf use is not advisable \pc{chanenson2025}{7}{this finding provides evidence that LLMs should not be applied off-the-shelf for GKC-CI annotation as of now}. Fine-tuned models reached 90.65\,\%, comparable to experts \pc{chanenson2025}{1}{Our best performing model has an accuracy of 90.65\%, which is com- parable to the accuracy of experts on the same task}. The open-source models tested did not converge within 10 epochs \pc{chanenson2025}{7}{No open-source model that we tested had performance high enough to be considered even a poor substi- tute}. Cost after training: \$0.23--0.44 per policy. The model was applied to 456 policies \pc{chanenson2025}{1}{We apply our best performing model to 456 privacy policies from a variety of online services}.

\cf{Error analysis} More than half of the apparent errors were not wrong: the model's span meant the same as the expert span, or the expert span was broader or narrower \pc{chanenson2025}{10}{We find that appreciably more than half of the purported match errors are the codes “Semantic Equivalence,” “Expansive Ground Truth,” and “Partial Ground Truth.”}. Exact-match scoring therefore understates span agreement.

\cf{Limits (card)} Models are of 2023 (GPT-3.5 fine-tuning, Llama 2); newer prompted models may do better, so the <20\,\% figure is an upper warning, not a forecast. Sixteen training policies from US web services. Fine-tuning needs thousands of labelled spans, which the thesis will not have. The authors' own limitation: policies are not written for the framework \pc{chanenson2025}{10}{website privacy policies are not written with the GKC-CI framework in mind, so a perfect mapping from text to GKC-CI parameters may not be possible}.

\cf{Use for this thesis} (1) Strongest caution for D9: theory-laden span coding (harm named, stance toward it, protection scope) may fail with prompting alone even where category detection works (Xie et al.\ \pp{xie2025}, Rodriguez et al.\ \pp{rodriguez2024}). The pilot must measure span-level recall per axis before the model screens anything. (2) The scoring rule matters: the methodology's ``right label, unit and evidence span'' is close to exact match; define in advance when an overlapping or semantically equivalent span counts, and report both strict and lenient scores. (3) CI parameters (sender, recipient, attribute, transmission principle, consequence) map onto the thesis's components and capabilities and onto Doan et al.\ \pp{doan2025}.
\end{papercard}
